% ============================================================================
% SPANISCH LANGENTWURF - DS 8: PUFFER-STUNDE
% ============================================================================
% Kurs: Jahrgangsübergreifend Spanisch Spätbeginner (Kl. 10, 11, 12)
% Datum: Di 27.01.2026
% Besonderheit: PUFFER - entbehrlich für Lehrprobe, Fokus auf Kreativität
% ============================================================================

\documentclass[11pt,a4paper]{article}

% ============================================================================
% PACKAGES
% ============================================================================
\usepackage[utf8]{inputenc}
\usepackage[T1]{fontenc}
\usepackage[ngerman]{babel}
\usepackage{geometry}
\usepackage{fancyhdr}
\usepackage{lastpage}
\usepackage[table]{xcolor}
\usepackage{tcolorbox}
\usepackage{enumitem}
\usepackage{tabularx}
\usepackage{longtable}
\usepackage{array}
\usepackage{booktabs}
\usepackage{hyperref}
\usepackage{ifthen}
\usepackage{amssymb}

% ============================================================================
% KONFIGURATION
% ============================================================================

% --- TIER ---
\newcommand{\plantier}{1}

% --- FACH ---
\newcommand{\fachid}{spanisch}
\newcommand{\fach}{Spanisch}
\newcommand{\fachfarbe}{c41e3a}

% ============================================================================
% VARIABLEN
% ============================================================================

% --- Metadaten ---
\newcommand{\jahrgangsstufe}{10/11/12 (jahrgangsübergreifend)}
\newcommand{\schulart}{Gesamtschule}
\newcommand{\stundenthema}{Creatividad y Juego -- Mi persona favorita}
\newcommand{\reihenthema}{Beobachtungsstunde Februar 2026}
\newcommand{\stundennummer}{8}
\newcommand{\reihenstunden}{8}
\newcommand{\unterrichtsdatum}{27.01.2026}
\newcommand{\unterrichtszeit}{[Lt. Stundenplan]}
\newcommand{\raum}{[Fachraum]}

% --- Lehrkraft ---
\newcommand{\lehrkraft}{[Name]}
\newcommand{\ausbildungsstand}{Lehrkraft}

% --- Lerngruppe ---
\newcommand{\lerngruppengroesse}{10}
\newcommand{\maennlich}{1}
\newcommand{\weiblich}{9}
\newcommand{\divers}{0}

% --- Kompetenzen/Niveau ---
\newcommand{\gerniveau}{A2--B1}
\newcommand{\lehrwerk}{A\_tope.com}
\newcommand{\lehrwerkseite}{--}

% ============================================================================
% LAYOUT
% ============================================================================
\geometry{left=2.5cm, right=2.5cm, top=2.5cm, bottom=2.5cm}
\definecolor{fach}{HTML}{\fachfarbe}
\definecolor{gray}{HTML}{666666}
\definecolor{lightgray}{HTML}{f5f5f5}
\definecolor{successgreen}{HTML}{2a7a4b}
\definecolor{warningorange}{HTML}{b35c00}
\definecolor{gruppeA}{HTML}{2c5aa0}
\definecolor{gruppeB}{HTML}{c41e3a}
\definecolor{puffer}{HTML}{6a0dad}
\definecolor{online}{HTML}{2a7a4b}

\tcbuselibrary{skins,breakable}

% Farbige Boxen
\newtcolorbox{infobox}[1][]{
    colback=lightgray,
    colframe=fach,
    fonttitle=\bfseries,
    title=#1,
    breakable
}

\newtcolorbox{scorebox}{
    colback=white,
    colframe=fach,
    fonttitle=\bfseries,
    title=Didaktik-Score,
    breakable
}

\newtcolorbox{pufferbox}{
    colback=white,
    colframe=puffer,
    fonttitle=\bfseries,
    title={PUFFER-STUNDE -- Entbehrlich für Lehrprobe},
    breakable
}

\newtcolorbox{onlinebox}{
    colback=white,
    colframe=online,
    fonttitle=\bfseries,
    title={ONLINE-VARIANTE -- Bei Lehrerausfall},
    breakable
}

% Kopf-/Fußzeile
\pagestyle{fancy}
\fancyhf{}
\fancyhead[L]{\small\textcolor{gray}{\fach{} -- Jg. \jahrgangsstufe{} -- \stundenthema}}
\fancyhead[R]{\small\textcolor{gray}{Langentwurf Tier 1}}
\fancyfoot[C]{\small\thepage{} / \pageref{LastPage}}
\renewcommand{\headrulewidth}{0.4pt}
\renewcommand{\footrulewidth}{0pt}

% ============================================================================
% DOKUMENT
% ============================================================================
\begin{document}

% ============================================================================
% DECKBLATT
% ============================================================================
\begin{titlepage}
    \centering
    \vspace*{2cm}

    {\large\textcolor{gray}{\schulart}}\\[2cm]

    {\Huge\textcolor{fach}{\textbf{Langentwurf}}}\\[0.5cm]
    {\large Tier 1 -- Vollständig}\\[2cm]

    {\LARGE\textbf{\stundenthema}}\\[0.5cm]
    {\large im Rahmen der Unterrichtsreihe}\\[0.3cm]
    {\Large\textit{\reihenthema}}\\[1cm]

    \textcolor{puffer}{\Large\textbf{--- PUFFER-STUNDE ---}}\\[0.3cm]
    {\normalsize Entbehrlich für Lehrprobe am 03.02.2026}\\[0.3cm]
    {\normalsize Fokus: Kreativität und Festigung ohne Prüfungsdruck}\\[1.5cm]

    \begin{tabular}{rl}
        \textbf{Fach:} & \fach{} (\gerniveau) \\[0.3cm]
        \textbf{Klasse:} & \jahrgangsstufe \\[0.3cm]
        \textbf{Lehrwerk:} & \lehrwerk \\[0.3cm]
        \textbf{Datum:} & \underline{\hspace{1cm}\unterrichtsdatum\hspace{1cm}} \\[0.3cm]
        \textbf{Zeit:} & \underline{\hspace{3cm}} (Doppelstunde) \\[0.3cm]
        \textbf{Raum:} & \underline{\hspace{3cm}} \\[0.3cm]
    \end{tabular}

    \vfill

    {\large Lehrkraft: \lehrkraft}\\[0.3cm]
    {\normalsize\textcolor{gray}{\ausbildungsstand}}\\[1cm]

\end{titlepage}

% ============================================================================
% INHALTSVERZEICHNIS
% ============================================================================
\tableofcontents
\newpage

% ============================================================================
% 1. BEDINGUNGSANALYSE
% ============================================================================
\section{Bedingungsanalyse}

\begin{pufferbox}
\textbf{Charakter dieser Stunde:} Diese Doppelstunde ist als \textbf{PUFFER} konzipiert und baut inhaltlich \textbf{nicht} auf die Beobachtungsstunde (DS 7, 03.02.2026) auf. Sie dient der Festigung bekannter Strukturen durch kreative Anwendung.

\textbf{Konsequenz:} Bei Zeitproblemen in der Sequenz kann diese Stunde gekürzt oder als Hausaufgabe ausgelagert werden, ohne die Beobachtungsstunde zu gefährden.
\end{pufferbox}

\subsection{Lerngruppe}
\begin{tabular}{ll}
    Kursgröße: & \lerngruppengroesse{} SuS (\maennlich{} m / \weiblich{} w / \divers{} d) \\
    Jahrgangsstufe: & Kl. 10, 11, 12 (jahrgangsübergreifend) \\
    Sprachniveau: & \gerniveau{} (GER) -- heterogen \\
    Fremdsprache: & 3. Fremdsprache (Spätbeginner) \\
\end{tabular}

\subsection{Schülerprofile}

\begin{tabular}{|l|c|c|l|p{5cm}|}
\hline
\textbf{Name} & \textbf{Kl.} & \textbf{Gruppe} & \textbf{NP} & \textbf{Profil} \\
\hline
\textcolor{gruppeA}{E.H.} & 10 & \textcolor{gruppeA}{A} & 15 & TOP, Peer-Tutor, Franz. Hintergrund \\
\hline
\textcolor{gruppeA}{L.N.} & 10 & \textcolor{gruppeA}{A} & $\sim$13 & Proaktiv, etwas faul \\
\hline
\textcolor{gruppeA}{C.P.} & 10 & \textcolor{gruppeA}{A} & 12 & Gute Aussprache, ruhig \\
\hline
\textcolor{gruppeA}{V.S.} & 10 & \textcolor{gruppeA}{A} & $\sim$11 & Konstant, fleißig, ruhig \\
\hline
\textcolor{gruppeB}{L.K.} & 11 & \textcolor{gruppeB}{B} & 10 & Nachholer, manchmal abgelenkt \\
\hline
\textcolor{gruppeB}{H.P.} & 11 & \textcolor{gruppeB}{B} & 10 & Späteinsteiger, Grundlagenlücken \\
\hline
\textcolor{gruppeB}{A.A.} & 12 & \textcolor{gruppeB}{B} & $\sim$10 & Fleißig, unsicher \\
\hline
\textcolor{gruppeB}{A.S.} & 12 & \textcolor{gruppeB}{B} & 9 & Einziger Junge, unsicher \\
\hline
\end{tabular}

\vspace{1em}

\textbf{Gruppen-Charakteristik:}
\begin{itemize}
    \item \textcolor{gruppeA}{\textbf{Gruppe A (5 SuS):}} A2-Niveau, Lernjahr 1, können frei sprechen
    \item \textcolor{gruppeB}{\textbf{Gruppe B (5 SuS):}} A1--B1 gemischt, Lernjahr 2--3, heterogener
\end{itemize}

\subsection{Besonderheiten}
\begin{itemize}
    \item \textbf{Puffer-Charakter:} Keine neuen Inhalte, nur kreative Anwendung
    \item \textbf{Wettbewerbselemente:} Staffel und Voting motivieren
    \item \textbf{Gamification:} Escape-Room-Rätsel als Highlight
    \item \textbf{Heterogenität:} Gallery Walk ermöglicht individuelle Vertiefung
\end{itemize}

% ============================================================================
% 2. SACHANALYSE
% ============================================================================
\section{Sachanalyse}

\subsection{Sprachliche Strukturen (Wiederholung)}

Diese Stunde wiederholt und festigt bekannte Strukturen aus DS 1--7:

\begin{itemize}
    \item \textbf{Adjektive zur Personenbeschreibung:}
    \begin{itemize}
        \item \textit{simpático/-a, inteligente, trabajador/-a, creativo/-a, divertido/-a}
        \item \textit{amable, paciente, responsable, organizado/-a, tranquilo/-a}
    \end{itemize}

    \item \textbf{Grammatische Strukturen:}
    \begin{itemize}
        \item \textit{ser + Adjektiv}: \textit{Mi madre es muy amable.}
        \item \textit{tener + años}: \textit{Tiene 45 años.}
        \item \textit{Possessivpronomen}: \textit{mi, su, nuestro/-a}
        \item \textit{Konnektoren}: \textit{también, además, pero, porque}
    \end{itemize}

    \item \textbf{Wortfelder:}
    \begin{itemize}
        \item Familie: \textit{madre, padre, hermano/-a, abuelo/-a}
        \item Berufe: \textit{profesor/-a, médico/-a, artista}
        \item Hobbys: \textit{le gusta + Infinitiv}
    \end{itemize}
\end{itemize}

\subsection{Kreative Methoden}

\textbf{1. Wortschatz-Staffel:}
\begin{itemize}
    \item Kompetitives Element fördert Abruf unter Zeitdruck
    \item Teamarbeit stärkt Gruppenzusammenhalt
    \item Wiederholung durch aktive Produktion
\end{itemize}

\textbf{2. Poster-Gestaltung:}
\begin{itemize}
    \item Visuelle Darstellung fördert Kreativität
    \item Kombination von Text und Bild
    \item Bekannte Struktur (\glqq{}Mi persona favorita\grqq{}) reduziert kognitive Last
\end{itemize}

\textbf{3. Gallery Walk:}
\begin{itemize}
    \item Peer-Learning durch gegenseitiges Lesen
    \item Voting schafft niedrigschwellige Bewertung
    \item Bewegung im Raum nach längerer Sitzphase
\end{itemize}

\textbf{4. Escape-Room-Rätsel:}
\begin{itemize}
    \item Problemlösendes Denken auf Spanisch
    \item Gamification steigert Motivation
    \item Selbstständiges Arbeiten in Kleingruppen
\end{itemize}

% ============================================================================
% 3. DIDAKTISCHE ANALYSE
% ============================================================================
\section{Didaktische Analyse}

\subsection{Stellung in der Sequenz}

Dies ist die \textbf{8. und letzte Doppelstunde} der Beobachtungssequenz.

\vspace{0.5em}
\textbf{Sequenzübersicht:}
\begin{center}
\begin{tabular}{|c|l|l|}
\hline
\textbf{DS} & \textbf{Thema} & \textbf{Fokus} \\
\hline
1--5 & Getrennte Progression & A: LJ1 / B: LJ2 \\
\hline
6 & Brückenthema: Personas y características & Erster Overlap \\
\hline
7 & Beobachtungsstunde (03.02.2026) & Vollständige Integration \\
\hline
\rowcolor{lightgray} \textbf{8} & \textbf{Creatividad y Juego} & \textbf{PUFFER -- Festigung} \\
\hline
\end{tabular}
\end{center}

\subsection{Puffer-Funktion}

\textbf{Warum Puffer?}
\begin{itemize}
    \item DS 7 (Beobachtungsstunde) ist der Höhepunkt der Sequenz
    \item DS 8 findet \textit{nach} der Beobachtung statt
    \item Kein neuer Stoff, der für die Lehrprobe relevant wäre
    \item Bei Zeitproblemen kann DS 8 gekürzt/ausgelagert werden
\end{itemize}

\textbf{Mehrwert als Abschluss:}
\begin{itemize}
    \item Positive Erfahrung zum Sequenzende
    \item Kreativität ohne Prüfungsdruck
    \item Reflexion des Gelernten durch Anwendung
\end{itemize}

\subsection{Legitimation}

\textbf{Rahmenplan Spanisch MV:}
\begin{itemize}
    \item \textbf{Schreiben:} Kurze kreative Texte verfassen
    \item \textbf{Sprechen:} Ergebnisse präsentieren
    \item \textbf{Wortschatz:} Aktiver Abruf und Anwendung
    \item \textbf{Methodenkompetenz:} Kooperative Lernformen
\end{itemize}

\textbf{Begründung kreative Methoden:}
\begin{itemize}
    \item Fremdsprachenlernen profitiert von emotional positiven Erfahrungen
    \item Spielerische Elemente senken die Sprechangst
    \item Peer-Feedback fördert Selbstreflexion
\end{itemize}

% ============================================================================
% 4. STUNDENZIELE
% ============================================================================
\section{Stundenziele}

\subsection{Grobziel}
Die SuS \textbf{festigen ihren Wortschatz und ihre Schreibkompetenz}, indem sie bekannte Strukturen kreativ anwenden und ihre Ergebnisse präsentieren.

\subsection{Feinziele (operationalisiert)}

\begin{tabular}{|c|p{6.5cm}|c|c|}
\hline
\textbf{Nr.} & \textbf{Die SuS können...} & \textbf{AFB} & \textbf{Phase} \\
\hline
FZ1 & mindestens 10 Adjektive zur Personenbeschreibung aktiv abrufen & I & Staffel \\
\hline
FZ2 & eine Person mit vollständigen Sätzen und Konnektoren beschreiben & II & Poster \\
\hline
FZ3 & ein visuell ansprechendes Poster mit Text und Bild gestalten & II & Poster \\
\hline
FZ4 & die Poster der Mitschüler lesen und bewerten & I/II & Gallery Walk \\
\hline
FZ5 & sprachliche Rätsel auf Spanisch lösen & II/III & Escape Room \\
\hline
\end{tabular}

\vspace{1em}
\textbf{Hinweis:} Reduzierter Anspruch (AFB I--II dominant), da Puffer-Stunde ohne neuen Stoff.

% ============================================================================
% 5. METHODISCHE ANALYSE
% ============================================================================
\section{Methodische Analyse}

\subsection{Phasierung (Doppelstunde = 90 Min.)}

\begin{center}
\begin{tabular}{|c|l|c|l|}
\hline
\textbf{Zeit} & \textbf{Phase} & \textbf{Min.} & \textbf{Sozialform} \\
\hline
0--15 & Wortschatz-Staffel & 15 & Teams (Plenum) \\
\hline
15--45 & Mini-Poster: \glqq{}Mi persona favorita\grqq{} & 30 & EA/PA \\
\hline
45--50 & Kurze Pause & 5 & -- \\
\hline
50--75 & Gallery Walk + Voting & 25 & Plenum (Bewegung) \\
\hline
75--90 & Escape-Room-Rätsel & 15 & GA (2--3er) \\
\hline
\rowcolor{lightgray} 90 & Auflösung + Prämierung & 5 & Plenum \\
\hline
\end{tabular}
\end{center}

\subsection{Detailplanung der Phasen}

\textbf{Phase 1: Wortschatz-Staffel (15 Min.)}
\begin{itemize}
    \item 2 Teams (gemischt A+B)
    \item Tafel geteilt: \glqq{}Team 1\grqq{} vs. \glqq{}Team 2\grqq{}
    \item Kategorie: \glqq{}Adjektive zur Personenbeschreibung\grqq{}
    \item Ablauf: Einer schreibt, rennt zurück, nächster
    \item Wertung: Korrekte Wörter zählen (Rechtschreibung!)
    \item Tiebreaker: Wer hat die \glqq{}schwierigsten\grqq{} Wörter?
\end{itemize}

\textbf{Phase 2: Mini-Poster (30 Min.)}
\begin{itemize}
    \item Aufgabe: \glqq{}Gestalte ein Poster über deine Lieblingsperson\grqq{}
    \item Vorgabe: Mind. 5 Sätze, mind. 1 Bild/Zeichnung
    \item Struktur vorgeben:
    \begin{enumerate}
        \item Wer ist die Person? (\textit{Mi persona favorita es...})
        \item Wie ist sie? (\textit{Es muy... y también...})
        \item Warum magst du sie? (\textit{Me gusta porque...})
    \end{enumerate}
    \item Material: A4-Papier, Stifte, Zeitschriften zum Ausschneiden
\end{itemize}

\textbf{Phase 3: Gallery Walk + Voting (25 Min.)}
\begin{itemize}
    \item Poster werden aufgehängt (Wände, Fenster)
    \item SuS gehen herum, lesen alle Poster
    \item Jede/r hat 3 Klebepunkte für Favoriten
    \item Kriterien: Sprache + Kreativität + Gesamteindruck
    \item Am Ende: Auszählung, Sieger-Poster prämieren
\end{itemize}

\textbf{Phase 4: Escape-Room-Rätsel (15 Min.)}
\begin{itemize}
    \item 3 Rätsel, jeweils 1 Code-Ziffer
    \item Rätsel kombinieren Wortschatz + Grammatik
    \item Richtiger 3-stelliger Code = \glqq{}Escape\grqq{}
    \item Belohnung: Kleine Süßigkeit / Lob
\end{itemize}

\subsection{Medien und Material}

\begin{itemize}
    \item Tafel (für Staffel)
    \item A4-Papier (10x), Buntstifte, Filzstifte
    \item Zeitschriften zum Ausschneiden (optional)
    \item Klebepunkte (30 Stück, 3 Farben)
    \item Escape-Room-Rätselblätter (3 Stück)
    \item Code-Schloss / Schatzkiste (Belohnung)
\end{itemize}

% ============================================================================
% 6. VERLAUFSPLANUNG PRÄSENZ
% ============================================================================
\section{Verlaufsplanung (Präsenz-Variante)}

\textbf{Übersicht Doppelstunde (90 Min.)}

\begin{longtable}{|p{0.8cm}|p{1.8cm}|p{3.8cm}|p{3cm}|p{1.3cm}|p{2cm}|}
\hline
\textbf{Zeit} & \textbf{Phase} & \textbf{Lehrerhandeln} & \textbf{Schülerhandeln} & \textbf{SF} & \textbf{Medien} \\
\hline
\endhead

\multicolumn{6}{|c|}{\textbf{WORTSCHATZ-STAFFEL (15 Min.)}} \\
\hline

0--2 & Einstieg & Begrüßung, Ankündigung: \glqq{}Hoy: creatividad y juegos!\grqq{} & Zuhören, Vorfreude & Plenum & -- \\
\hline

2--4 & Team-bildung & Teams einteilen (gemischt A+B), Tafel teilen & In Teams aufstellen & Plenum & Tafel \\
\hline

4--12 & Staffel & Startsignal, Zeitnehmer, bei Bedarf korrigieren & Rennen zur Tafel, Adjektive schreiben & Teams & Tafel, Kreide \\
\hline

12--15 & Auswertung & Wörter zählen, Sieger küren, schwierigste Wörter loben & Jubeln, Ergebnisse vergleichen & Plenum & -- \\
\hline

\multicolumn{6}{|c|}{\textbf{MINI-POSTER (30 Min.)}} \\
\hline

15--18 & Einführung & Aufgabe erklären: \glqq{}Mi persona favorita\grqq{}, Struktur zeigen & Zuhören, Fragen stellen & Plenum & Tafel \\
\hline

18--20 & Material & Material verteilen (Papier, Stifte) & Material holen & -- & Material \\
\hline

20--45 & Arbeits-phase & Herumgehen, Hilfe anbieten, sprachlich korrigieren & Poster gestalten: schreiben, zeichnen & EA/PA & Papier, Stifte \\
\hline

\multicolumn{6}{|c|}{\textbf{PAUSE (5 Min.)}} \\
\hline

45--50 & Pause & Poster aufhängen lassen, Raum vorbereiten & Pause, Poster aufhängen & -- & Klebeband \\
\hline

\multicolumn{6}{|c|}{\textbf{GALLERY WALK + VOTING (25 Min.)}} \\
\hline

50--52 & Erklärung & Gallery Walk erklären, Klebepunkte verteilen (3 pro Person) & Punkte entgegennehmen & Plenum & Klebepunkte \\
\hline

52--65 & Gallery Walk & Beobachten, bei Bedarf Vokabeln erklären & Herumgehen, lesen, Punkte kleben & Bewegung & Poster \\
\hline

65--72 & Auszählung & Punkte zählen, Ranking erstellen & Zählen helfen, Spannung & Plenum & -- \\
\hline

72--75 & Prämierung & Sieger-Poster vorstellen, Applaus, evtl. Preis & Applaudieren, Sieger gratulieren & Plenum & -- \\
\hline

\multicolumn{6}{|c|}{\textbf{ESCAPE-ROOM-RÄTSEL (15 Min.)}} \\
\hline

75--77 & Einführung & Escape-Room erklären: 3 Rätsel = 3-stelliger Code & Zuhören, Gruppen bilden & Plenum & -- \\
\hline

77--88 & Rätselphase & Rätselblätter verteilen, bei Bedarf Tipps geben & Rätsel lösen in Kleingruppen & GA (3er) & Rätselblätter \\
\hline

88--90 & Auflösung & Code auflösen, Gewinner-Gruppe prämieren, Abschluss & Code eingeben, jubeln & Plenum & Code-Schloss \\
\hline

\end{longtable}

\textbf{Legende:} EA = Einzelarbeit, PA = Partnerarbeit, GA = Gruppenarbeit, SF = Sozialform

% ============================================================================
% 7. ONLINE-VARIANTE (BEI LEHRERAUSFALL)
% ============================================================================
\section{Online-Variante (bei Lehrerausfall)}

\begin{onlinebox}
Diese Variante ist für den Fall konzipiert, dass die Lehrkraft ausfällt und die SuS selbstständig arbeiten müssen. Alle Aufgaben sind über \textbf{itslearning} abrufbar und abzugeben.

\textbf{Zeitrahmen:} 90 Min. selbstständiges Arbeiten
\end{onlinebox}

\subsection{Aufgabe 1: Schreibaufgabe (40 Min.)}

\textbf{Thema:} \glqq{}Una persona importante en mi vida\grqq{}

\textbf{Aufgabenstellung (auf itslearning):}
\begin{quote}
Escribe un texto sobre una persona importante en tu vida. Puede ser un familiar, un/-a amigo/-a, un/-a profesor/-a o una persona famosa.

\textbf{Tu texto debe incluir:}
\begin{itemize}
    \item Quién es la persona (nombre, relación contigo)
    \item Cómo es (mínimo 5 adjetivos)
    \item Por qué es importante para ti
    \item Un ejemplo / una anécdota
\end{itemize}

\textbf{Extensión:} 80--120 palabras

\textbf{Estructura recomendada:}
\begin{enumerate}
    \item Introducción: \textit{Una persona muy importante en mi vida es...}
    \item Descripción: \textit{Es muy... y también... Además,...}
    \item Razón: \textit{Es importante para mí porque...}
    \item Ejemplo: \textit{Por ejemplo, una vez...}
\end{enumerate}
\end{quote}

\textbf{Abgabe:} Textfeld oder Datei-Upload in itslearning

\subsection{Aufgabe 2: Selbsttest (25 Min.)}

\textbf{Format:} Google Form oder itslearning-Quiz

\textbf{Inhalt (20 Fragen):}
\begin{enumerate}
    \item 10 Fragen: Adjektiv-Zuordnung (Deutsch -- Spanisch)
    \item 5 Fragen: \textit{ser} vs. \textit{estar} (Lückentext)
    \item 3 Fragen: Possessivpronomen ergänzen
    \item 2 Fragen: Sätze ordnen
\end{enumerate}

\textbf{Automatische Auswertung:} Sofortiges Feedback nach Abgabe

\subsection{Aufgabe 3: Padlet-Beitrag (optional, 15 Min.)}

\textbf{Falls Zeit übrig:}
\begin{itemize}
    \item Padlet-Link in itslearning
    \item Aufgabe: Poste ein Bild + 3 Sätze über \glqq{}Mi persona favorita\grqq{}
    \item Kommentiere mindestens 2 Beiträge von Mitschülern
\end{itemize}

\subsection{Übersicht Online-Variante}

\begin{center}
\begin{tabular}{|c|l|c|l|}
\hline
\textbf{Nr.} & \textbf{Aufgabe} & \textbf{Zeit} & \textbf{Abgabe} \\
\hline
1 & Schreibaufgabe & 40 Min. & itslearning \\
\hline
2 & Selbsttest & 25 Min. & Google Form / itslearning \\
\hline
3 & Padlet (optional) & 15 Min. & Padlet \\
\hline
\end{tabular}
\end{center}

\textbf{Hinweis für Vertretungslehrkraft:} Die SuS loggen sich in itslearning ein und arbeiten die Aufgaben selbstständig ab. Die Lehrkraft kann bei technischen Problemen helfen, muss aber keine fachliche Unterstützung geben.

% ============================================================================
% 8. ERWARTETE SCHWIERIGKEITEN
% ============================================================================
\section{Erwartete Schwierigkeiten}

\begin{tabular}{|p{4.5cm}|p{8.5cm}|}
\hline
\textbf{Schwierigkeit} & \textbf{Reaktion} \\
\hline
\textbf{Staffel: Nur wenige Adjektive bekannt} & Vorab kurze Wiederholung (2 Min.): \glqq{}Welche Adjektive kennen wir?\grqq{} als Warm-up \\
\hline
\textbf{Staffel: Ein Team deutlich schwächer} & Teams nachträglich ausgleichen oder Handicap-Regel (schwächeres Team startet früher) \\
\hline
\textbf{Poster: SuS schreiben zu wenig} & Mindestanforderung visualisieren (5 Sätze), Struktur-Vorlage an die Tafel \\
\hline
\textbf{Poster: SuS schreiben auf Deutsch} & Nur spanische Texte werden aufgehängt. Bei Bedarf Wörterbuch-App erlauben. \\
\hline
\textbf{Gallery Walk: Alle Punkte für Freunde} & Regel: Nicht das eigene Poster bewerten. Mind. 1 Punkt für Sprachqualität vergeben. \\
\hline
\textbf{Escape Room: Zu schwer} & Gestufte Hilfen vorbereiten (Umschläge mit Tipps). L gibt nach 5 Min. ersten Hinweis. \\
\hline
\textbf{Escape Room: Zu schnell gelöst} & Bonus-Rätsel vorbereiten oder schnelle Gruppen helfen langsameren \\
\hline
\textbf{Online: SuS arbeiten nicht} & Abgabe als Pflicht kommunizieren, Zeitstempel prüfbar. Nicht-Abgabe = 0 Punkte (falls benotet) \\
\hline
\end{tabular}

% ============================================================================
% 9. ANLAGEN
% ============================================================================
\section{Anlagen}

\begin{enumerate}
    \item Poster-Vorlage: \glqq{}Mi persona favorita\grqq{}
    \item Escape-Room-Rätsel (3 Blätter)
    \item Bewertungskriterien Gallery Walk
    \item Online-Aufgabenstellung (für itslearning)
    \item Selbsttest-Fragen (für Google Form)
\end{enumerate}

\subsection{Anlage 1: Poster-Vorlage}

\begin{verbatim}
+---------------------------------------------------------------------+
|                     MI PERSONA FAVORITA                              |
+---------------------------------------------------------------------+
|                                                                     |
|  [PLATZ FÜR BILD / ZEICHNUNG]                                       |
|                                                                     |
+---------------------------------------------------------------------+
|                                                                     |
|  Mi persona favorita es _________________________________           |
|                                                                     |
|  Es mi __________________________ (madre, amigo, profesor...)       |
|                                                                     |
|  Es muy __________________ y también __________________.            |
|                                                                     |
|  Además, es __________________ pero no es __________________.       |
|                                                                     |
|  Me gusta mucho porque ______________________________________       |
|                                                                     |
|  _____________________________________________________________      |
|                                                                     |
+---------------------------------------------------------------------+
\end{verbatim}

\subsection{Anlage 2: Escape-Room-Rätsel}

\textbf{Rätsel 1: Wortschlange (Code-Ziffer 1)}

\begin{verbatim}
Finde 6 versteckte Adjektive in der Schlange:

SIMPÁTICOINTELIGENTECREATIVOPACIENTERESPONSABLEDIVERTIDO

Schreibe die Anfangsbuchstaben hintereinander: _ _ _ _ _ _

Die 3. Buchstabe (als Zahl: A=1, B=2, C=3...) ist die erste
Code-Ziffer.

Lösung: S I C P R D -> C = 3 -> Code-Ziffer 1: 3
\end{verbatim}

\textbf{Rätsel 2: Lückentext (Code-Ziffer 2)}

\begin{verbatim}
Ergänze ser oder estar:

1. Mi hermano _____ muy simpático. (ser)
2. Hoy yo _____ cansada. (estar)
3. Ella _____ enferma. (estar)
4. Nosotros _____ estudiantes. (ser)
5. ¿Tú _____ contento? (estar)

Zähle: Wie oft steht "ser"? ___
Das ist die zweite Code-Ziffer.

Lösung: ser kommt 2x vor -> Code-Ziffer 2: 2
\end{verbatim}

\textbf{Rätsel 3: Zuordnung (Code-Ziffer 3)}

\begin{verbatim}
Ordne die Adjektive den Berufen zu:

A) médico/-a      1) creativo/-a
B) profesor/-a    2) valiente
C) artista        3) paciente
D) bombero/-a     4) organizado/-a

Richtige Zuordnung: A-3, B-4, C-1, D-2

Addiere die Zahlen der richtigen Zuordnungen für
A und B: 3 + 4 = ?

Das ist die dritte Code-Ziffer.

Lösung: 3 + 4 = 7 -> Code-Ziffer 3: 7
\end{verbatim}

\textbf{Gesamtcode: 3 - 2 - 7}

\subsection{Anlage 3: Bewertungskriterien Gallery Walk}

\begin{verbatim}
+---------------------------------------------------------------------+
|  KRITERIEN FÜR DIE BEWERTUNG (Gallery Walk)                         |
+---------------------------------------------------------------------+
|                                                                     |
|  Vergib deine 3 Punkte nach diesen Kriterien:                       |
|                                                                     |
|  1. SPRACHE                                                         |
|     - Sind die Sätze auf Spanisch?                                  |
|     - Sind sie grammatisch korrekt?                                 |
|     - Werden verschiedene Adjektive verwendet?                      |
|                                                                     |
|  2. KREATIVITÄT                                                     |
|     - Ist das Poster schön gestaltet?                               |
|     - Gibt es ein Bild oder eine Zeichnung?                         |
|     - Ist es originell?                                             |
|                                                                     |
|  3. GESAMTEINDRUCK                                                  |
|     - Möchte ich mehr über diese Person erfahren?                   |
|     - Ist das Poster interessant?                                   |
|                                                                     |
|  REGELN:                                                            |
|  - Du darfst NICHT dein eigenes Poster bewerten                     |
|  - Du kannst mehrere Punkte auf ein Poster kleben                   |
|  - Oder die Punkte auf verschiedene Poster verteilen                |
|                                                                     |
+---------------------------------------------------------------------+
\end{verbatim}

% ============================================================================
% 10. DIDAKTIK-SCORE
% ============================================================================
\newpage
\section{Didaktik-Score}

\begin{scorebox}
\textbf{Bewertung der didaktischen Qualität nach 10 Dimensionen}\\
\textbf{Ziel-Score: $\geq$ 9.0 (Premium-Qualität)}

\vspace{1em}
\begin{tabular}{|c|l|c|c|p{4.5cm}|}
\hline
\textbf{\#} & \textbf{Dimension} & \textbf{Gew.} & \textbf{Score} & \textbf{Notiz} \\
\hline
1 & Differenzierung & 15\% & 8/10 & Implizit durch offene Aufgaben, nicht explizit gesteuert \\
\hline
2 & Taxonomie (AFB) & 10\% & 8/10 & AFB I--II dominant (Puffer), wenig AFB III \\
\hline
3 & Lernziel-Alignment & 10\% & 10/10 & Festigung bekannter Strukturen, RP-konform \\
\hline
4 & Aktivierung & 10\% & 10/10 & Staffel, Poster, Gallery Walk = hohe Aktivierung \\
\hline
5 & Scaffolding & 10\% & 9/10 & Vorlagen, Struktur, gestufte Hilfen bei Escape Room \\
\hline
6 & Feedback-Qualität & 10\% & 9/10 & Peer-Feedback (Voting), L korrigiert bei Poster \\
\hline
7 & Sprachliche Klarheit & 10\% & 10/10 & Einfache Aufgabenstellungen, bekannte Strukturen \\
\hline
8 & Motivation \& Relevanz & 10\% & 10/10 & Spielerisch, Wettbewerb, keine Noten \\
\hline
9 & Inklusion (UDL) & 10\% & 9/10 & Visuell + kinästhetisch, Bewegung, aber auditiv wenig \\
\hline
10 & Praktikabilität & 5\% & 10/10 & Material einfach, Zeit realistisch, Online-Alternative \\
\hline
\multicolumn{3}{|r|}{\textbf{GESAMT:}} & \textbf{9.2/10} & \\
\hline
\end{tabular}

\vspace{1em}
\textbf{Bewertungsschwellen:}
\begin{itemize}
    \item[\checkmark] \textcolor{successgreen}{$\geq$ 9.0: \textbf{FREIGABE} (Premium-Qualität)}
    \item[$\Box$] \textcolor{warningorange}{8.0--8.9: Iteration empfohlen}
    \item[$\Box$] 6.0--7.9: Überarbeitung erforderlich
    \item[$\Box$] $<$ 6.0: Grundlegende Neukonzeption
\end{itemize}

\vspace{1em}
\textcolor{successgreen}{\textbf{STATUS: FREIGABE}}

\vspace{1em}
\textbf{Stärken:}
\begin{itemize}
    \item Hohe Motivation durch spielerische Elemente (Staffel, Escape Room)
    \item Kreative Freiheit ohne Prüfungsdruck
    \item Online-Variante ermöglicht Flexibilität bei Ausfall
    \item Peer-Feedback durch Gallery Walk
    \item Bewegung im Raum nach längeren Sitzphasen
\end{itemize}

\textbf{Bewusste Einschränkungen (Puffer-Charakter):}
\begin{itemize}
    \item Keine explizite Niveaudifferenzierung (offene Aufgaben stattdessen)
    \item Wenig AFB III (bewusst: Festigung, nicht Neulernen)
    \item Kein neuer Stoff (Sequenz-Abschluss)
\end{itemize}

\end{scorebox}

\end{document}
